{{دیوان
| عنوان = بچّہ اور شمع
| ماخذ = ویکی ماخذ
| ماخذ_لنک = https://ur.wikisource.org/wiki/%D8%A8%DA%86%DB%81%D9%91_%D8%A7%D9%88%D8%B1_%D8%B4%D9%85%D8%B9
| مقابلہ = اقبال سائبر لائبریری (اقبال اکادمی پاکستان)
| مقابلہ_لنک = https://www.iqbalcyberlibrary.net
}}
<poem>
کیسی حیرانی ہے یہ اے طفلکِ پروانہ خُو!
شمع کے شعلوں کو گھڑیوں دیکھتا رہتا ہے تو

یہ مری آغوش میں بیٹھے ہوئے جُنبش ہے کیا
روشنی سے کیا بغَل گیری ہے تیرا مدّعا؟

اس نظارے سے ترا ننھّا سا دل حیران ہے
یہ کسی دیکھی ہُوئی شے کی مگر پہچان ہے

شمع اک شعلہ ہے لیکن تُو سراپا نور ہے
آہ! اس محفل میں یہ عُریاں ہے تُو مستور ہے

دستِ قُدرت نے اسے کیا جانے کیوں عُریاں کیا!
تُجھ کو خاکِ تِیرہ کے فانوس میں پنہاں کیا

نور تیرا چھُپ گیا زیرِ نقابِ آگہی
ہے غبارِ دیدۂ بینا حجابِ آگہی

زندگانی جس کو کہتے ہیں فراموشی ہے یہ
خواب ہے، غفلت ہے، سرمستی ہے، بے ہوشی ہے یہ

محفلِ قُدرت ہے اک دریائے بے پایانِ حُسن
آنکھ اگر دیکھے تو ہر قطرے میں ہے طوفانِ حُسن

حُسن، کوہستاں کی ہیبت ناک خاموشی میں ہے
مِہر کی ضوگستری، شب کی سِیہ پوشی میں ہے

آسمانِ صبح کی آئینہ پوشی میں ہے یہ
شام کی ظُلمت، شفَق کی گُل فرو شی میں ہے یہ

عظمتِ دیرینہ کے مِٹتے ہُوئے آثار میں
طفلکِ ناآشنا کی کوششِ گُفتار میں

ساکنانِ صحنِ گُلشن کی ہم آوازی میں ہے
ننھّے ننھّے طائروں کی آشیاں سازی میں ہے

چشمۂ کُہسار میں، دریا کی آزادی میں حُسن
شہر میں، صحرا میں، ویرانے میں، آبادی میں حُسن

رُوح کو لیکن کسی گُم گشتہ شے کی ہے ہوس
ورنہ اس صحرا میں کیوں نالاں ہے یہ مثلِ جرس!

حُسن کے اس عام جلوے میں بھی یہ بے تاب ہے
زندگی اس کی مثالِ ماہیِ بے آب ہے
</poem>
